\section{Class Reference}
\label{sec:class-reference}

This appendix describes all 20 classes defined in the \MLIPsOntology{}.

\subsection{\dlConcept{AccuracyMetric}}
\label{sec:term-AccuracyMetric}

\textbf{Label:} Accuracy Metric.
An accuracy measure reported in a benchmark
(e.g., RMSE of energy, MAE of forces).

\paragraph{Outgoing properties.} $\dlRole{metricProperty}$, $\dlRole{metricType}$, $\dlRole{metricValue}$.

\paragraph{Incoming properties.} $\dlRole{hasAccuracyMetric}$.

\paragraph{Example.}
\InputIfFileExists{artifacts/examples/classes/AccuracyMetric.tex}{}{\textit{TODO: add example.}}

\paragraph{Related axioms.} (A24), (A25), (A27).

\subsection{\dlConcept{AtomicConfiguration}}
\label{sec:term-AtomicConfiguration}

\textbf{Label:} Atomic Configuration.
A single atomic configuration (positions, cell,
properties) in a training dataset.

\paragraph{Incoming properties.} $\dlRole{hasConfiguration}$.

\paragraph{Example.}
\InputIfFileExists{artifacts/examples/classes/AtomicConfiguration.tex}{}{\textit{TODO: add example.}}

\subsection{\dlConcept{BenchmarkResult}}
\label{sec:term-BenchmarkResult}

\textbf{Label:} Benchmark Result.
A single evaluation result: an algorithm applied to a
material system with specific hyperparameters, reporting accuracy metrics.

\paragraph{Outgoing properties.} $\dlRole{evaluatesModel}$, $\dlRole{hasAccuracyMetric}$, $\dlRole{hasHyperparameterSetting}$, $\dlRole{reportedIn}$, $\dlRole{targetMaterial}$, $\dlRole{usesAlgorithm}$, $\dlRole{usesTrainingData}$.

\paragraph{Incoming properties.} $\dlRole{hasResult}$.

\paragraph{Example.}
\InputIfFileExists{artifacts/examples/classes/BenchmarkResult.tex}{}{\textit{TODO: add example.}}

\paragraph{Related axioms.} (A21), (A22), (A23), (A26).

\subsection{\dlConcept{BenchmarkStudy}}
\label{sec:term-BenchmarkStudy}

\textbf{Label:} Benchmark Study.
A published study reporting MLIP evaluation results.

\paragraph{Superclasses.} $\dlConcept{Activity}$.

\paragraph{Outgoing properties.} $\dlRole{hasResult}$.

\paragraph{Example.}
\InputIfFileExists{artifacts/examples/classes/BenchmarkStudy.tex}{}{\textit{TODO: add example.}}

\paragraph{Related axioms.} (A20).

\subsection{\dlConcept{CoveredProperty}}
\label{sec:term-CoveredProperty}

\textbf{Label:} Covered Property.
A physical property covered by the training dataset
(energy, forces, stresses, virials).

\paragraph{Incoming properties.} $\dlRole{coversProperty}$.

\paragraph{Example.}
\InputIfFileExists{artifacts/examples/classes/CoveredProperty.tex}{}{\textit{TODO: add example.}}

\subsection{\dlConcept{DFTCalculation}}
\label{sec:term-DFTCalculation}

\textbf{Label:} DFT Calculation.
A density functional theory calculation that produced
reference data for training.

\paragraph{Outgoing properties.} $\dlRole{hasDFTSettings}$.

\paragraph{Incoming properties.} $\dlRole{hasDFTCalculation}$.

\paragraph{Example.}
\InputIfFileExists{artifacts/examples/classes/DFTCalculation.tex}{}{\textit{TODO: add example.}}

\paragraph{Related axioms.} (A19).

\subsection{\dlConcept{DFTSettings}}
\label{sec:term-DFTSettings}

\textbf{Label:} DFT Settings.
Settings of a DFT calculation: exchange-correlation
functional, k-point mesh, energy cutoff, pseudopotential type.

\paragraph{Outgoing properties.} $\dlRole{usedDFTCode}$, $\dlRole{energyCutoff}$, $\dlRole{kPointMesh}$, $\dlRole{pseudopotentialType}$, $\dlRole{xcFunctional}$.

\paragraph{Incoming properties.} $\dlRole{hasDFTSettings}$.

\paragraph{Example.}
\InputIfFileExists{artifacts/examples/classes/DFTSettings.tex}{}{\textit{TODO: add example.}}

\subsection{\dlConcept{DatasetProvenance}}
\label{sec:term-DatasetProvenance}

\textbf{Label:} Dataset Provenance.
The provenance type of a training dataset.

\paragraph{Incoming properties.} $\dlRole{datasetProvenance}$.

\paragraph{Example.}
\InputIfFileExists{artifacts/examples/classes/DatasetProvenance.tex}{}{\textit{TODO: add example.}}

\subsection{\dlConcept{Hyperparameter}}
\label{sec:term-Hyperparameter}

\textbf{Label:} Hyperparameter.
A hyperparameter of an MLIP algorithm, with its name,
data type, range, and default value.

\paragraph{Superclasses.} $\dlConcept{HyperParameter}$.

\paragraph{Outgoing properties.} $\dlRole{defaultValue}$, $\dlRole{hyperparameterDatatype}$, $\dlRole{hyperparameterName}$, $\dlRole{maxValue}$, $\dlRole{minValue}$.

\paragraph{Incoming properties.} $\dlRole{forHyperparameter}$, $\dlRole{hasHyperparameter}$.

\paragraph{Example.}
\InputIfFileExists{artifacts/examples/classes/Hyperparameter.tex}{}{\textit{TODO: add example.}}

\paragraph{Related axioms.} (A8).

\subsection{\dlConcept{HyperparameterSetting}}
\label{sec:term-HyperparameterSetting}

\textbf{Label:} Hyperparameter Setting.
A concrete value assigned to a hyperparameter in a
specific training run or configuration.

\paragraph{Outgoing properties.} $\dlRole{forHyperparameter}$, $\dlRole{cutoffRadius}$, $\dlRole{learningRate}$, $\dlRole{numAngularBasis}$, $\dlRole{numLayers}$, $\dlRole{numRadialBasis}$, $\dlRole{settingValue}$.

\paragraph{Incoming properties.} $\dlRole{hasHyperparameterSetting}$.

\paragraph{Example.}
\InputIfFileExists{artifacts/examples/classes/HyperparameterSetting.tex}{}{\textit{TODO: add example.}}

\subsection{\dlConcept{Implementation}}
\label{sec:term-Implementation}

\textbf{Label:} Implementation.
A software implementation of an MLIP algorithm in a
specific library and version.

\paragraph{Outgoing properties.} $\dlRole{implementedIn}$, $\dlRole{version}$.

\paragraph{Incoming properties.} $\dlRole{hasImplementation}$.

\paragraph{Example.}
\InputIfFileExists{artifacts/examples/classes/Implementation.tex}{}{\textit{TODO: add example.}}

\paragraph{Related axioms.} (A3), (A9).

\subsection{\dlConcept{Library}}
\label{sec:term-Library}

\textbf{Label:} Library.
A software library or tool used in MLIP research
(e.g., ASE, LAMMPS, VASP, GPAW).

\paragraph{Incoming properties.} $\dlRole{implementedIn}$, $\dlRole{usedDFTCode}$.

\paragraph{Example.}
\InputIfFileExists{artifacts/examples/classes/Library.tex}{}{\textit{TODO: add example.}}

\subsection{\dlConcept{MLIPAlgorithm}}
\label{sec:term-MLIPAlgorithm}

\textbf{Label:} MLIP Algorithm.
A machine learning interatomic potential algorithm family
(e.g., MACE, ACE, HDNNP, NequIP, M3GNet, MTP).

\paragraph{Superclasses.} $\dlConcept{Algorithm}$.

\paragraph{Outgoing properties.} $\dlRole{hasHyperparameter}$, $\dlRole{hasImplementation}$, $\dlRole{supportsSimulation}$.

\paragraph{Incoming properties.} $\dlRole{executes}$, $\dlRole{trainedWith}$, $\dlRole{usesAlgorithm}$.

\paragraph{Example.}
\InputIfFileExists{artifacts/examples/classes/MLIPAlgorithm.tex}{}{\textit{TODO: add example.}}

\paragraph{Related axioms.} (A1), (A2).

\subsection{\dlConcept{MaterialSystem}}
\label{sec:term-MaterialSystem}

\textbf{Label:} Material System.
A material system studied in a dataset: element, alloy,
compound, or complex microstructure.

\paragraph{Outgoing properties.} $\dlRole{chemicalFormula}$, $\dlRole{materialClass}$, $\dlRole{microstructuralFeature}$.

\paragraph{Incoming properties.} $\dlRole{coversMaterial}$, $\dlRole{targetMaterial}$.

\paragraph{Example.}
\InputIfFileExists{artifacts/examples/classes/MaterialSystem.tex}{}{\textit{TODO: add example.}}

\subsection{\dlConcept{MetricProperty}}
\label{sec:term-MetricProperty}

\textbf{Label:} Metric Property.
The physical property being measured (energy, force, stress).

\paragraph{Incoming properties.} $\dlRole{metricProperty}$.

\paragraph{Example.}
\InputIfFileExists{artifacts/examples/classes/MetricProperty.tex}{}{\textit{TODO: add example.}}

\subsection{\dlConcept{MetricType}}
\label{sec:term-MetricType}

\textbf{Label:} Metric Type.
Type of accuracy metric (RMSE, MAE, R², etc.).

\paragraph{Incoming properties.} $\dlRole{metricType}$.

\paragraph{Example.}
\InputIfFileExists{artifacts/examples/classes/MetricType.tex}{}{\textit{TODO: add example.}}

\subsection{\dlConcept{SimulationType}}
\label{sec:term-SimulationType}

\textbf{Label:} Simulation Type.
A type of atomistic simulation: molecular dynamics,
Monte Carlo, geometry optimization, phonon calculations,
thermodynamic integration, etc.

\paragraph{Incoming properties.} $\dlRole{supportsSimulation}$.

\paragraph{Example.}
\InputIfFileExists{artifacts/examples/classes/SimulationType.tex}{}{\textit{TODO: add example.}}

\subsection{\dlConcept{TrainedModel}}
\label{sec:term-TrainedModel}

\textbf{Label:} Trained Model.
The concrete artifact produced by training an algorithm
on a specific dataset with specific hyperparameter settings---the learned
parameters (weights) that can be used for prediction.

\paragraph{Superclasses.} $\dlConcept{Model}$.

\paragraph{Outgoing properties.} $\dlRole{trainedOn}$, $\dlRole{trainedWith}$.

\paragraph{Incoming properties.} $\dlRole{evaluatesModel}$, $\dlRole{produces}$.

\paragraph{Example.}
\InputIfFileExists{artifacts/examples/classes/TrainedModel.tex}{}{\textit{TODO: add example.}}

\paragraph{Related axioms.} (A7).

\subsection{\dlConcept{TrainingDataset}}
\label{sec:term-TrainingDataset}

\textbf{Label:} Training Dataset.
A dataset of atomic configurations used for MLIP
training, with provenance, size, property coverage, and DFT settings.

\paragraph{Superclasses.} $\dlConcept{Dataset}$, $\dlConcept{Entity}$.

\paragraph{Outgoing properties.} $\dlRole{coversMaterial}$, $\dlRole{coversProperty}$, $\dlRole{datasetProvenance}$, $\dlRole{hasConfiguration}$, $\dlRole{hasDFTCalculation}$, $\dlRole{numConfigurations}$.

\paragraph{Incoming properties.} $\dlRole{runsOn}$, $\dlRole{trainedOn}$, $\dlRole{usesTrainingData}$.

\paragraph{Example.}
\InputIfFileExists{artifacts/examples/classes/TrainingDataset.tex}{}{\textit{TODO: add example.}}

\paragraph{Related axioms.} (A15), (A16), (A17), (A18).

\subsection{\dlConcept{TrainingRun}}
\label{sec:term-TrainingRun}

\textbf{Label:} Training Run.
A concrete training execution: it executes an algorithm
on a training dataset and produces a trained model.

\paragraph{Superclasses.} $\dlConcept{Run}$, $\dlConcept{Activity}$.

\paragraph{Outgoing properties.} $\dlRole{executes}$, $\dlRole{produces}$, $\dlRole{runsOn}$.

\paragraph{Example.}
\InputIfFileExists{artifacts/examples/classes/TrainingRun.tex}{}{\textit{TODO: add example.}}

\paragraph{Related axioms.} (A4), (A5), (A6).
